%% EXEMPLO FICTÍCIO — Atividade 6: evento, competição ou visita técnica
\begin{atividade}{Competição internacional de ciclismo com estimulação elétrica funcional}{
  tipo         = evento,
  modalidade   = EVT,
  alternativa  = COMP,
  horas        = 60,
  inicio       = 2024-10,
  fim          = 2024-10,
  projeto      = {NEBULA, estimulação elétrica funcional para ciclismo},
  papel        = {Técnica responsável pelo estimulador},
  supervisor   = {Rafael Fictício, supervisor do projeto NEBULA},
  registros    = {EXT-901; NRO-DIR-006-901},
  competencias = {IV, V, VI},
}

\begin{contexto}
A competição reúne equipes de vários países em provas com tecnologias assistivas; na prova de ciclismo
com estimulação elétrica funcional, o atleta pedala uma distância fixa, e o sistema é avaliado também
por uma banca técnica. Para o NEBULA, é o teste mais duro que existe: o sistema tem de funcionar longe
do laboratório, depois de uma viagem, sob o regulamento de outra instituição.
\end{contexto}

\begin{contribuicao}
Fui a técnica responsável pelo estimulador: montei o checklist de preparação, levei uma placa reserva
configurada, fiz o plano de carga das baterias, operei o estimulador nas provas e apresentei a parte
eletrônica do sistema à banca técnica, em inglês.
\end{contribuicao}

\begin{desenvolvimento}
\fichatecnica{
  evento     = {Competição internacional de tecnologias assistivas, prova de ciclismo com FES},
  local      = {Europa},
  datas      = {Quatro dias em outubro de 2024},
  funcao     = {Técnica responsável pelo estimulador},
  resultado  = {4.º lugar entre 9 equipes},
  declaracao = {\texttt{NRO-DIR-006-901}},
}

\textbf{A preparação} (20 horas, em setembro). O checklist de preparação cobriu a placa reserva, os
eletrodos, os cabos e o transporte das baterias de lítio, que as companhias aéreas só aceitam abaixo de
\SI{100}{\watt\hour} por unidade e com a documentação do fabricante --- as nossas têm
\SI{29}{\watt\hour}. Simulei, com o atleta e a fisioterapeuta, duas provas completas no LABBIO.

\textbf{A competição} (quatro dias de cerca de dez horas). Conferência dos equipamentos na chegada,
treinos livres, a apresentação à banca técnica e as provas. Em cada prova, eu conferia a impedância dos
eletrodos, carregava o perfil de estímulo e acompanhava a telemetria.
\end{desenvolvimento}

\begin{resultados}
A equipe terminou em 4.º lugar entre 9. O estimulador funcionou em todas as provas; um eletrodo descolou
na segunda, e o limite de corrente por hardware desligou o canal como devia --- o atleta completou a
prova com três canais. Da conversa com as outras equipes, trouxemos a ideia do controle da cadência em
malha fechada, que virou requisito do NEBULA para 2025.
\end{resultados}

\begin{evidencias}
\evidencia{Declaração}{NRO-DIR-006-901}{A declaração de participação, com a função e as horas}{auth.neurodynamics.dev (código verificador)}
\evidencia{Registro}{EXT-901}{O registro do evento no SOMA, aprovado por duas pessoas}{SOMA › Serviços › Eventos}
\evidencia{Resultado}{Classificação oficial}{A colocação da equipe na prova}{Site da competição}
\end{evidencias}

\begin{aprendizados}
Operar um sistema meu diante de uma banca de outro país ensinou mais sobre preparação do que qualquer
teste de bancada: o checklist salvou a segunda prova. E ver sistemas de outras equipes mostrou o que o
nosso ainda não fazia.
\end{aprendizados}
\end{atividade}
